% Write last. Three points only: parameter count, timing, takeaway.
% Accuracy: at most one sentence saying it is not the subject of this paper.
Hybrid quantum--classical physics-informed neural networks (QCPINNs) are usually reported to be
efficient since they have fewer trainable parameters than a classical PINN. We ask how much
wall-clock time this saving costs. We evaluate a classical PINN together with four hybrid
configurations that vary in qubit count and circuit depth on a one-dimensional seawater
temperature diffusion benchmark.
\emph{Parameter count:} the quantum branch adds only 16--64 trainable parameters, and the
hybrids have 2421--2717 parameters in total against 4353 for the classical baseline. This amounts
to a reduction of 38--44\%, almost all of it due to the smaller classical portion.
\emph{Timing:} relative to the classical baseline, the four hybrid configurations are
14--59$\times$ slower on the forward pass, 14--102$\times$ slower for the first-order gradient,
and 13--94$\times$ slower for the second-order gradient, growing with qubit count, circuit depth
and batch size.
\emph{Takeaway:} the number of parameters alone tends to overstate the efficiency of hybrid
models. This paper is not about accuracy.
